We separate note-, report-, inventory-, helper-, and mirror-style benign cases into a dedicated hard-benign pack because they are the dominant false-positive pressure point for explanation-oriented runtime security analysis. They are outward-looking by presentation and often overlap with risky command or artifact patterns, yet they remain benign under the benchmark definition because they terminate locally, use public or low-sensitivity content, or send content only to expected sinks. Treating this slice as first-class evaluation material prevents benchmark optimism from being driven only by malicious-case coverage and makes residual false positives directly inspectable.
